\subsection{换位子}

定义 3. 设 G 为一个群，x 与 y 为 G 的两个元素。G 的元素 \(x^{-1}y^{-1}xy\) 称为 x 与 y 的换位子。

用 \((x,y)\) 表示 x 与 y 的换位子。于是显然

\[
(y, x) = (x, y) ^ {- 1}.
\]

对于 \(x\) 与 \(y\) ，可交换的充要条件是 \((x, y) = e\) 。更一般地，

\[
x y = y x (x, y).
\]

另一方面，我们记

\[
x ^ {y} = y ^ {- 1} x y = x (x, y) = (y, x ^ {- 1}) x.\tag{3}
\]

由于映射 \(x \mapsto x^y\) 是内自同构 \(\operatorname{Int}(y^{-1})\) ，故 \((x, y)^z = (x^z, y^z)\) 对所有 \(x, y, z \in G\) 成立。

对于 \(x, y, z \in G\) ，我们证明以下关系式：

\[
(x, y z) = (x, z). (x, y) ^ {z} = (x, z). (z, (y, x)). (x, y)\tag{4}
\]

(4 bis) \((xy,z) = (x,z)^y.(y,z) = (x,z).(((x,z),y).(y,z))\)

\begin{enumerate}
\def\labelenumi{(\arabic{enumi})}
\setcounter{enumi}{4}
\item
  \((x^{y}, (y, z)).(y^{z}, (z, x)).(z^{x}, (x, y)) = e\)
\item
  \((x, yz).(z, xy).(y, zx) = e\)
\end{enumerate}

(6 bis) \((xy, z). (yz, x). (zx, y) = e.\)

现在

\[
\begin{array}{r l} (x, y z) & = x ^ {- 1} z ^ {- 1} y ^ {- 1} x y z = (x, z) z ^ {- 1} x ^ {- 1} y ^ {- 1} x y z = (x, z) (x, y) ^ {z} \\ & = (x, z) (z, (x, y) ^ {- 1}) (x, y) \end{array}
\]

由 (3) 即得，这证明了 (4)。公式 (4 bis) 可类似地得到。另一方面，

\[
\begin{array}{r l} (x ^ {y}, (y, z)) & = (x ^ {y}) ^ {- 1} (z, y) (x ^ {y}) (y, z) \\ & = y ^ {- 1} x ^ {- 1} y z ^ {- 1} y ^ {- 1} z y y ^ {- 1} x y y ^ {- 1} z ^ {- 1} y z \\ & = (y z y ^ {- 1} x y) ^ {- 1} (z x z ^ {- 1} y z). \end{array}
\]

然后记 \(u = yzy^{-1}xy\) , \(v = zxz^{-1}yz\) 和 \(w = xyx^{-1}zx\) ，我们得到

\[
(x ^ {y}, (y, z)) = u ^ {- 1} v.
\]

通过循环置换 \(x, y, z\) ，我们推出 \((y^z, (z, x)) = v^{-1}w\) 和

\[
\left(z ^ {x}, (x, y)\right) = w ^ {- 1} u,
\]

这立即推出(5)。最后，(6)通过将x, y, z在公式 \((x, yz) = x^{-1}z^{-1}y^{-1}xyz = (yzx)^{-1}(xyz)\) 中循环置换所得三个公式的两边分别相乘而得到，对于(6 bis)也类似。

若 A 和 B 是 G 的两个子群，则 (A, B) 表示由诸换位子 \((a, b)\) （ \(a \in A\) ， \(b \in B\) ）生成的子群。于是 \((A, B) = \{e\}\) 当且仅当 A 中心化 B； \((A, B) \subset A\) 当且仅当 B 正规化 A。若 A 和 B 都是正规的（相应地，特征的），则 \((A, B)\) 亦然。

命题 5. 设 A、B、C 是 G 的三个子群。

\begin{enumerate}
\def\labelenumi{(\roman{enumi})}
\item
  子群A正规化子群(A, B)。
\item
  若子群 (B, C) 正规化 A，则子群 (A, (B, C)) 由诸元素 \((a, (b, c))\) （ \(a \in \mathbf{A}\) , \(b \in \mathbf{B}\) , \(c \in \mathbf{C}\) ）生成。
\item
  若 A、B 和 C 均为正规子群，则
\end{enumerate}

\[
(\mathrm{A}, (\mathrm{B}, \mathrm{C})) \subset (\mathrm{C}, (\mathrm{B}, \mathrm{A})). (\mathrm{B}, (\mathrm{C}, \mathrm{A})).
\]

由 (4 bis)，对 \(a, a' \in \mathbf{A}\) 与 \(b \in \mathbf{B}\) ，

\[
(a, b) ^ {a ^ {\prime}} = (a a ^ {\prime}, b). (a ^ {\prime}, b) ^ {- 1},
\]

由此得 (i)。现在假设 (B, C) 正规化 A。对于 \(a \in \mathbf{A}\) , \(b \in \mathbf{B}\) , \(c \in \mathbf{C}\) 和 \(x \in \mathbf{G}\) ，(4) 蕴含

\[
(a, (b, c). x) = (a, x). (x, ((b, c), a)) (a, (b, c))
\]

和 \(((b, c), a) \in \mathbf{A}\) ，因为 (B, C) 正规化 A；故由对 \(p\) 的归纳可知， \(\left(a, \prod_{i=1}^{p} (b_i, c_i)\right)\) （对 \(b_i \in \mathbf{B}\) , \(c_i \in \mathbf{C}\) ）属于由形如 \((a, (b, c))\) 的元素生成的子群。若最后 A、B、C 都是正规的，则子群 (A, (B, C))、(C, (B, A)) 与 (B, (C, A)) 也都是正规的。于是由 (ii)，只需证明

\[
(a, (b, c)) \in (\mathrm{C}, (\mathrm{B}, \mathrm{A})). (\mathrm{B}, (\mathrm{C}, \mathrm{A}))
\]

对所有 \(a \in \mathbf{A}\) , \(b \in \mathbf{B}\) 与 \(c \in \mathbf{C}\) 成立。现在由 (5)，记 \(a^{b^{-1}} = u\)

\[
(a, (b, c)) = ((u ^ {b}), (b, c)) = (c ^ {u}, (u, b)) ^ {- 1}. (b ^ {c}, (c, u)) ^ {- 1},
\]

由此得(iii)。

定义 4. 设 G 是一个群。由 G 中元素的换位子生成的子群称为 G 的导群。

G 的导群即子群 (G, G)。它也记作 D(G)。作为一种语言上的滥用，它有时被称为 G 的换位子群，尽管它一般不同于 G 中元素的换位子集合（习题 16）。D(G) = \{e\} 当且仅当 G 是交换的。

命题 6. 设 \(f: \mathbf{G} \to \mathbf{G}'\) 是群同态。则 \(f(\mathbf{D}(\mathbf{G})) \subset \mathbf{D}(\mathbf{G}')\) 。若 \(f\) 是满射，则 \(\mathbf{D}(\mathbf{G})\) 到 \(\mathbf{D}(\mathbf{G}')\) 的、由 \(f\) 的限制所得到的同态是满射。

\(f\) 把 \(G\) 中元素的换位子映为 \(G'\) 中元素的换位子。若 \(f\) 是满射，则 \(f\) 把 \(G\) 的换位子集合映为 \(G'\) 的换位子集合。于是该命题由 §4 第 3 号命题 2 的推论 3 得出。

推论1. 群G的导群是G的一个特征子群。特别地，它是G的一个正规子群。

推论 2. 设 G 是一个群。商群 G/D(G) 是交换的。设 \(\pi: G \to G/D(G)\) 为典范同态。每个同态 \(f\) （从 G 到交换群 \(G'\) ）都可唯一地表示为 \(f = \bar{f} \circ \pi\) 的形式，其中 \(\bar{f}: G/D(G) \to G'\) 是一个同态。

现在， \(\pi (\mathbf{D}(\mathbf{G})) = \{e\}\) 。由于 \(\pi\) 是满射，由此可得 \(\mathbf{D}(\mathbf{G} / \mathbf{D}(\mathbf{G})) = \{e\}\) ，从而得出第一个断言。第二个断言由 §4 第 4 号命题 5 得出。

推论 3. 设 H 是 G 的一个子群。下列条件等价：

\begin{enumerate}
\def\labelenumi{(\roman{enumi})}
\item
  \(\mathbf{H} \supset \mathbf{D}(\mathbf{G})\) ;
\item
  H是正规子群且G/H是交换群。
\item
  (ii) \(\Rightarrow\) (i) 由推论 2 可得，而 (i) \(\Rightarrow\) (ii) 由 §4 第 7 号定理 4 可得，因为交换群的每个子群都是正规的。
\end{enumerate}

推论 4. 设 G 是一个群，X 是 G 的一个生成 G 的子集。群 D(G) 是由 X 中元素的换位子生成的 G 的正规子群。

设 H 是由 X 中元素的换位子生成的 G 的正规子群，且 \(\phi: G \to G/H\) 为典范同态。集合 \(\phi(X)\) 生成 G/H。\(\phi(X)\) 中的元素两两可交换，因此 H 是交换的（ \(\S 4\) ，第 3 号，命题 2 的推论 2）。于是（推论 3）H 包含 D(G)。另一方面，显然 \(H \subset D(G)\) 。

注记。(1) 推论2也可以表述为：G/D(G)连同 \(\pi\) ，是G关于交换群及从G到交换群的同态的泛映射问题的解。

\begin{enumerate}
\def\labelenumi{(\arabic{enumi})}
\setcounter{enumi}{1}
\tightlist
\item
  在推论4的假设下，由X中元素的换位子生成的子群包含于D(G)中，但一般不等于D(G)（参见习题15e）。
\end{enumerate}

例子。(1) 若 \(\mathbf{G}\) 是非交换单群，则 \(\mathrm{D}(\mathbf{G}) = \mathbf{G}\) 。因此，从 \(\mathbf{G}\) 到交换群的每个同态都是平凡的。

\begin{enumerate}
\def\labelenumi{(\arabic{enumi})}
\setcounter{enumi}{1}
\tightlist
\item
  对称群 \(\mathfrak{S}_n\) 的导群是交错群 \(\mathfrak{A}_n\) 。因为 \(\mathfrak{A}_n\) 由两个对换的乘积生成；若 \(\tau = \tau_{x,y}\) 与 \(\tau' = \tau_{x',y'}\) 是两个对换，设 \(\sigma\) 为满足 \(\sigma(x') = x\) 与 \(\sigma(y') = y\) 的置换。则 \(\tau' = \sigma^{-1}\tau\sigma\) ，且 \(\tau\tau' = \tau^{-1}\tau' = \tau^{-1}\sigma^{-1}\tau\sigma\) 是一个换位子。因此 \(\mathfrak{A}_n \subset D(\mathfrak{S}_n)\) 。由于 \(\mathfrak{S}_n / \mathfrak{A}_n\) 是交换的，故 \(\mathfrak{A}_n \supset D(\mathfrak{S}_n)\) （推论 3）。
\end{enumerate}

